\section{Evidence and Supported Conclusions}
\label{sec:evidence}

% Role: setup. Message: separate the three evidence strata.
We evaluate the framework through three distinct evidence strata: historical adapted model studies, bounded live harness-role pilots, and source-bound engineering qualification. The historical inventory contains fourteen adapted method profiles and one Native mechanism profile across sixteen recipes; one profile can contain more than one recipe. These are archived experiments from earlier source revisions, not reruns of the current release-preparation tree. The pilot matrix tests specified role bindings on small tasks. Engineering checks test execution contracts and recovery rather than model quality. Appendix~\ref{app:evidence} links each stratum to its source and records what was inspected for this report.

% Role: setup and boundary. Message: retain original model and budget semantics.
Historical AIME and Formula configurations use the recorded \code{deepseek-flash} model identifier with temperature zero and thinking disabled. AIME uses an 8,192-token response limit; the Formula roles use 4,096. These are the identifiers and settings in the archived configurations, not assertions of equivalence to an upstream paper's model or to an unreported provider-resolved revision. Model seeds do not establish deterministic remote inference. We retain baseline, selected, and last-candidate distinctions, report task counts directly, and avoid pooled cross-task rankings or single-seed significance claims.

\subsection{Historical task results: examples of supported study protocols}
\label{sec:historical}
% Role: evidence. Message: different evaluation protocols require different labels.
The AIME and standalone Formula studies illustrate why a shared record format must retain evaluation semantics. GEPA/AIME uses 45 training, 45 validation, and 30 private tasks; its recorded baseline solves 23 private tasks and its selected candidate solves 24. ACE/Formula uses 500/300/200 tasks. Offline ACE records 141/200 for the initial playbook and 173/200 for the selected playbook, while online ACE records 163/200 pre-learning stream predictions against a 141/200 baseline. The online number describes predictions made before each learning window, not an untouched test of the final playbook. Table~\ref{tab:historical} keeps these quantities separate.

\begin{table}[t]
\caption{Historical adapted single-seed results. Counts are correct tasks / evaluated tasks. Offline selected-candidate results and online stream predictions have different meanings; these rows are not a cross-task ranking.}
\label{tab:historical}
\small
\centering
\begin{tabular}{@{}lrrl@{}}
\toprule
Profile & Initial & Reported outcome & Evaluation meaning \\
\midrule
GEPA / AIME & 23/30 & 24/30 & Selected, private \\
ACE / Formula offline & 141/200 & 173/200 & Selected, private \\
ACE / Formula online & 141/200 & 163/200 & Pre-learning stream \\
\bottomrule
\end{tabular}
\end{table}


% Role: interpretation. Message: selection provenance matters even with complete results.
Offline ACE also shows why selecting the last artifact would change the reported experiment. The retained candidate is selected after 200 training examples, whereas the last playbook follows all 500. Only the initial and selected playbooks have private measurements; the last playbook's private score is unavailable. Proteus records the selected candidate, the final execution artifact, their hashes, and the visible-only selection rule separately. The 16-percentage-point observed difference is a result of this one adapted trajectory, not a claim about expected gain or the last candidate.

% Role: evidence. Message: code artifacts enable a distinct adapted protocol.
The historical Meta-Harness-inspired study covers USPTO, Symptom2Disease, and LawBench with NoMemory and FewShotAll baselines plus three generated memory programs. Its proposer makes one JSON call that returns three programs, rather than conducting the original paper's extended tool-using coding search. Each system has 450 learning callbacks and 260 visible predictions over 130 distinct visible tasks. Four selected systems receive 412 private tasks each. Table~\ref{tab:mce} reports their per-domain counts; the generated systems do not dominate FewShotAll. This experiment demonstrates a code-artifact and learning/evaluation protocol under its recorded adaptation, while current coding-proposer support is assessed separately by pilots.

\begin{table}[t]
\caption{Historical Meta-Harness-inspired MCE private results. The proposer is one JSON call producing three programs; four systems were selected for private evaluation. The unselected lesson-distilled program has no private result. Counts are correct / tasks, and generated systems do not dominate FewShotAll.}
\label{tab:mce}
\small
\centering
\begin{tabular}{@{}llrrr@{}}
\toprule
System & Origin & USPTO & Symptom2Disease & LawBench \\
\midrule
NoMemory & Baseline & 16/100 & 151/212 & 28/100 \\
FewShotAll & Baseline & 29/100 & 183/212 & 55/100 \\
Contrastive boundary & Generated & 29/100 & 162/212 & 45/100 \\
Dual-pool retrieval & Generated & 30/100 & 152/212 & 52/100 \\
\bottomrule
\end{tabular}
\end{table}


% Role: evidence and limitation. Message: multiworker completion is not cooperation benefit.
The historical CORAL-inspired circle-packing arms compare one proposer making 20 submissions with two proposers making ten each, using one evaluation worker. The recorded selected normalized scores are 0.928722 and 0.900208, respectively, from an initial 0.364102. Twelve versus eleven submissions succeed; eight versus nine time out. Final confirmation repeats the same public 26-circle problem. These results do not establish a multiworker advantage, hidden-task generalization, or the behavior of the original persistent native agents. The locally unrecovered historical cross-worker chains are not repaired by relabeling a later protocol check.

\subsection{Composition: historical measurements and a controlled trace}
\label{sec:mechanism}
% Role: evidence. Message: costs and timings have their own observable boundaries.
The Combee/Formula inventory contains sequential ACE, naive parallel, fixed Combee, and dynamic Combee in both offline and online modes. The four online conditions have accepted pre-learning accuracy, independently recorded API upper-cost estimates, and study elapsed measurements (Table~\ref{tab:combee}). The API amounts are not invoices and do not overwrite unknown canonical native USD totals. Study elapsed time includes the recorded study boundary and differs from the sum of overlapping operation durations. Shared host resources and continuation histories prevent interpreting the table as an isolated speedup experiment; the different baseline scores also remain visible.

\begin{table}[t]
\caption{Historical online Combee/Formula conditions, seed 0. Accuracy is pre-learning stream correct / 200. API upper cost is USD, not an invoice or canonical native cost. Elapsed hours follow the recorded study timer; shared resources and continuations preclude an isolated speedup claim.}
\label{tab:combee}
\small
\centering
\begin{tabular}{@{}lrrrr@{}}
\toprule
Condition & Initial & Stream & API upper (USD) & Study elapsed (h) \\
\midrule
ACE & 139/200 & 160/200 & 1.738 & 4.62 \\
Naive & 137/200 & 157/200 & 0.705 & 2.81 \\
Fixed & 140/200 & 159/200 & 0.772 & 3.00 \\
Dynamic & 137/200 & 162/200 & 0.742 & 2.89 \\
\bottomrule
\end{tabular}
\end{table}


% Role: evidence. Message: fixed inputs isolate composition mechanics.
A separate model-free check feeds eleven frozen original task/prediction/feedback/reflection records to the four production ACE variants. All arms use the same H0 and seed. Sequential ACE curates each example; naive parallel uses batches 3/3/3/2; fixed Combee uses the same batch schedule with copied groups and reduction; dynamic Combee uses 1/2/4/4 under a controlled clock. Curator invocation counts are 11, 4, 12, and 12, respectively, while every arm makes 22 generator invocations. Provider-call count is zero. The auditor binds actual payload bytes, source/configuration locks, grouping, requests, updates, and lineage. This establishes distinct, reusable composition semantics under common inputs; synthetic timing and deterministic responses cannot establish learned quality or model speedup.

\subsection{Harness-role interoperability: sixteen bounded live cells}
\label{sec:pilots}
% Role: evidence. Message: explicit cells support role-level portability.
The live pilot inventory contains sixteen accepted cells across Pi, DSH, and Codex (Table~\ref{tab:pilots}). Native and GEPA cells use actual Pi/DSH source adapters as task subjects; GEPA's reflector is a separately managed chat actor. Other cells put the coding harness in the ACE/Combee generator role, Meta-Harness proposer role, or a bounded CORAL worker wrapper. The tasks are deliberately small and their role placement is part of the evidence. A checkmark therefore means that the named minimum cycle was observed, not that all roles of a complete benchmark ran on that harness.

\begin{table}[t]
\caption{Sixteen accepted bounded live cells. \yes{} = the listed minimum cycle has source-bound live evidence; \na{} = not established by this pilot inventory. The harness role differs by row. These are integration pilots, not full benchmarks or comparisons of model quality.}
\label{tab:pilots}
\small
\begin{tabularx}{\linewidth}{@{}p{0.18\linewidth}cccY@{}}
\toprule
Method/profile & Pi & DSH & Codex & Role exercised by the named harness \\
\midrule
Native common & \yes & \yes & \na & Evolution actor and candidate subject; one episode. \\
GEPA & \yes & \yes & \na & Candidate subject; separate managed chat reflection. \\
ACE & \yes & \yes & \yes & Generator executing the current playbook. \\
Combee-ACE & \yes & \yes & \yes & Independent generator workspaces; one map/reduce batch. \\
Meta-Harness & \yes & \yes & \yes & Coding proposer; generated memory-system programs execute. \\
CORAL bounded & \yes & \yes & \yes & One supervised worker; two fresh coding workspaces. \\
Combee-GEPA & \na & \na & \na & Implemented composition, outside this live matrix. \\
\bottomrule
\end{tabularx}
\end{table}


% Role: interpretation. Message: successful integration can include no quality gain.
Successful cycles include outcomes that do not improve on their baselines. The Pi ACE and Combee-ACE pilots modify and execute a playbook but retain the initial candidate, with private scores unchanged at zero. Meta-Harness coding proposers generate executable systems, yet the selected baselines tie or exceed them on the observed pilot tasks; those generated systems need not have private scores. Other cells demonstrate a changed policy or source artifact, independent scoring, and zero additional dispatches on completed resume. These outcomes support interoperability and auditability while resisting an artificial all-positive performance story.

% Role: boundary. Message: live cells do not prove universal compatibility.
The matrix leaves Codex Native/GEPA and all Combee-GEPA live role cells unclaimed. The additional CORAL role pilots use one supervised worker invoking fresh coding workspaces, not persistent native-session multiworker execution. Requested/applied model identifiers are preserved separately from missing provider-resolved identities. Cell-level evidence spans several recorded commits and includes corrected failed attempts. It is appropriate for the bounded integration claim, but cannot be assembled into a single-source release gate or a controlled comparison of harness quality.

\subsection{Current engineering and directed protocol qualification}
% Role: evidence. Message: direct delivery is stronger than availability alone.
A directed CORAL qualification specifically tests shared-content delivery, intervention, and controller-crash recovery. The accepted live run contains exactly three completed paid provider requests, reporting 1,141 input tokens and 113 output tokens; USD cost is unknown. Readback follows the actual delivered content, applied intervention, original request identities, and settled receipts. Resume performs the declared final confirmation evaluation without an additional model request. This closes a narrow protocol chain that historical knowledge snapshots alone could not prove. It does not establish worker-crash recovery, native model-session continuation, simultaneous provider dispatch, or a benefit from collaboration.

% Role: evidence. Message: a full preparation smoke is mechanism qualification.
The complete preparation smoke, run \href{https://github.com/proteus-evolve/Proteus/actions/runs/41398194611}{41398194611}, passes Offline, Minimal, LLM, Pi, DSH, and Codex jobs in one source-bound run. The offline job reports 1,842 passing tests and 91 skips; the live jobs include the declared source-edit, activation, and rejection checks. This is engineering qualification. Pi/DSH diagnostic archives are bounded and sanitized: Pi's 20 phase outcomes and 20 of DSH's 24 outcomes return 137 under controller caps, with incomplete usage streams. Codex's original strict-checker log establishes its checked edit/binary behavior, but its raw trace, snapshot, and binary archive were not exported. Passing this bounded gate neither proves complete billing nor makes the retained artifacts a full replay archive.

% Role: evidence. Message: reports can be reconstructed without new inference.
Independent generic readback reconstructs nine report/CSV files for each of five saved runs, covering 92 task rows and 30 candidates with no unreported rollout. Costs and timing intervals that were missing remain missing. Installation evidence covers 22 checks outside the checkout, including extension templates, exports, and study comparison. These results support the usability of the records and public extension boundary. They do not measure researcher productivity, and the current report does not rerun models, containers, or benchmarks to reproduce them.

% Role: interpretation. Message: each evidence stratum answers a different question.
Taken together, the evidence shows implemented algorithm profiles, concrete role-level execution, and inspectable state transitions under declared conditions. Historical task improvements illustrate that the substrate can carry substantive studies; negative pilot outcomes and retained failures show why implementation acceptance is a separate criterion. Controlled trace checks expose reusable mechanics, and directed qualification tests a specific recovery chain. None of these alone establishes a general self-improvement advantage, a universal harness interface, or a successful future release tag.
